\chapter{Conclusion}

\section{What the Thesis Presents}

The thesis presents a system that starts from real League of Legends match data and from it builds player profiles, relationships, graphs, clusters, analytic results, and finally a simulation experiment. Beyond kill and death counts, the system records which players a given player ends up in a game with, who they face as opponents, what recurring relationships form around them, and what larger network patterns emerge from those.

If the essence of the project had to be reduced to a single point, it would be the existence of the full chain rather than any individual component. A Dijkstra implementation or a React page in isolation is a routine project. What makes the system interesting is the connectedness: from raw match data, through the graph, clusters, and Rust analytics, it reaches a reproducible simulation experiment, and every step is verifiable.

\section{Main Results}

The data-collection layer handles two different queue types in parallel. Keeping \textit{flexset} and \textit{soloq} separate is important because the two queues do not measure the same phenomenon: Flex Queue is better suited for examining recurring connection patterns; SoloQ serves as a control and an individual performance baseline.

The \textit{opscore} and \textit{feedscore} metrics provide role-aware scores. They are not perfect truth measurements, but they allow performance indicators to be used as inputs to graph analytics.

The \textit{flexset} graph has a core-periphery structure: there are denser, more strongly connected core regions and many looser, smaller peripheral groups. This follows from the queue mechanics.

The assortativity analysis shows a positive association for \textit{opscore} in social-path mode (0.235): connected players carry more similar performance profiles than would be expected under random connections.

The Rust-based algorithmic layer --- pathfinding, assortativity, betweenness centrality --- is deterministic and reproducible: the same input data always produces the same result.

\section{The Role of Tribal NeuroSim}

Tribal NeuroSim is a transfer experiment: it examines whether the performance profiles extracted from the graph and its clusters can be carried over into the initial parameters of an evolutionary, multi-agent simulation.

Based on the four runs, profiles from the Flex Queue showed more consistent behaviour: a similar Warband/Combat winning archetype emerged under both tested seeds. The SoloQ runs were more variable --- different seeds produced different winning profiles. The flexset priors thus produced more consistent outcomes; the soloq priors produced more variable ones --- under an identical simulation rule set.

\section{Retrospective: What Worked Well}

In retrospect, the best decision was not forcing everything into a single technology. Python stayed for rapid data collection and experimentation; Rust took over where performance and determinism were needed; Node/Express and React made the whole thing presentable. Forcing all of this into a single language would have made at least one layer suffer.

The project's strength is the coherent chain: from raw API data through normalised player profiles, graph construction, and clustering, all the way to analytic results --- every step is verifiable and re-runnable. Assortativity, betweenness centrality, and the core-periphery interpretation are not just visualisation; they ask real questions about the network.

\section{Limitations and Open Questions}

A few limitations are worth keeping in mind when interpreting the results.

The dataset is not a random sample. The seed player selection, the premade-probe mechanism, and the queue type together determine which player population enters the graph. The resulting structure is a product of the given collection strategy, not a general description of all EUNE Master+ players.

opscore and feedscore are useful and interpretable metrics, but not external ground-truth measures. The role expectations and curve parameters are part of the model --- different construction choices would yield different results.

Even in the graph, an enemy connection is not equivalent to genuine conflict. Repeated opposition can arise from matchmaking MMR bands, shared time-of-day, or the fact that the active Master+ Flex Queue player pool is inherently small. This is why the signed graph remained an experimental option.

Temporal consistency and community cohesion analyses fell out of the current scope. The excluded analyses are not uninteresting; the data and methodology did not, however, support defensible claims.

\section{Closing Remarks}

The end result of the project can be stated simply: it is possible to build a system from competitive game data that holds up from a software engineering perspective and can also pose research questions. PremadeGraph is a data-collection, data-modelling, graph-analytic, and simulation chain.

The connectivity between steps was not planned in advance. Raw match data became normalised player profiles, profiles became a relationship graph, the graph became clusters and analytic metrics, and clusters became simulation priors --- each step taking the previous step's output as its input.

As ABBA put it decades ago: \textit{the winner takes it all}.
In Tribal NeuroSim, this is not a metaphor but a direct consequence of the mechanics.

\section{Where This Applies Elsewhere}
\label{sec:osszegzes-transfer}

The methodology is not tied to League of Legends. The Riot API was a convenient choice because it is publicly available, rich, and structured. But the same pipeline would run equally well on GitHub collaboration data, NBA performance statistics, or PubMed co-authorship networks --- anywhere there is recurring interaction, measurable performance, and some form of group structure.

Graph clustering, for example, describes a corporate organisational network just as naturally as a player population: in both cases the question is which actual working groups emerge behind the formal structure. Betweenness centrality measures the same thing in infrastructure planning or epidemiology --- which node's removal most disrupts the network.

The five Tribal NeuroSim artifacts are not LoL-specific concepts either. Combat, resource, map\_objective, risk, team --- these translate to other domains, and the evolutionary simulation framework can start from any cluster profile where group-level performance indicators can be extracted. The video game here was a tool, not a constraint.
